\notechapter{N-20230205}{Vector Triple Product, Exterior Algebra, and the Hodge Star}

\section{The vector triple product}

The note begins with a proof of the identity
\[
  u\times(v\times w)=(u\cdot w)v-(u\cdot v)w.
\]
The key observation is that \(u\times(v\times w)\) lies in the plane spanned by \(v\) and \(w\), because it is perpendicular to \(v\times w\).  Therefore
\[
  u\times(v\times w)=A v+B w.
\]
Taking dot products with suitable vectors determines \(A\) and \(B\), giving
\[
  A=u\cdot w,\qquad B=-(u\cdot v).
\]

This proof is better than expanding a determinant because it explains the geometry: the result must live in the \(v,w\)-plane, and only the two scalar projections of \(u\) onto \(v,w\) can appear.

\section{Scalar triple product}

The scalar triple product is
\[
  [u,v,w]=(u\times v)\cdot w.
\]
It equals the signed volume of the parallelepiped spanned by \(u,v,w\).  It is alternating: swapping two vectors changes the sign.  It vanishes exactly when the three vectors are coplanar.

The note uses this to justify identities involving cross products and dot products.  The recurring geometric idea is that determinants measure oriented volume.

\section{From vector identities to exterior algebra}

This note should not try to become a general tensor-calculus note.  Its natural center is narrower and cleaner: the familiar three-dimensional vector identities are shadows of the metric, the volume form, wedge product, contraction, and the Hodge star.

\begin{definition}[Bivectors and the wedge product]
For vectors \(u,v\) in a vector space \(V\), the wedge product \(u\wedge v\) is an alternating two-vector.  It records the oriented area element spanned by \(u\) and \(v\), not a third vector perpendicular to them.
\end{definition}

In a three-dimensional oriented Euclidean space, the metric and orientation identify bivectors with vectors.  This identification is the Hodge star:
\[
  \star:\Lambda^2 V\longrightarrow V.
\]
The cross product is then
\[
  u\times v=\star(u\wedge v).
\]
Thus the cross product is not the primitive operation.  The primitive operation is the oriented area element \(u\wedge v\); the cross product appears only after choosing both a metric and an orientation.

\begin{remark}[Why dimension three is special]
In dimension \(3\), bivectors and vectors both have dimension \(3\), so the Hodge star can identify them.  In other dimensions, \(\Lambda^2 V\) and \(V\) usually have different dimensions, and there is no ordinary vector-valued cross product with the same properties.
\end{remark}


\section{Levi-Civita notation}

The coordinate proof becomes much shorter with the Levi-Civita symbol.  Write
\[
  (u\times v)_i=\varepsilon_{ijk}u_jv_k,
\]
using the Einstein summation convention.  Then
\[
  (u\times(v\times w))_i
  =\varepsilon_{ijk}u_j\varepsilon_{klm}v_lw_m.
\]
The contraction identity
\[
  \varepsilon_{ijk}\varepsilon_{klm}
  =\delta_{il}\delta_{jm}-\delta_{im}\delta_{jl}
\]
gives
\[
  (u\times(v\times w))_i=(u_jw_j)v_i-(u_jv_j)w_i.
\]
This is exactly
\[
  u\times(v\times w)=(u\cdot w)v-(u\cdot v)w.
\]

This notation is useful because it separates the combinatorics of signs from the geometry of the identity.

\section{Lagrange identity}

Another important identity is
\[
  (u\times v)\cdot(w\times z)
  =(u\cdot w)(v\cdot z)-(u\cdot z)(v\cdot w).
\]
Setting \(w=u\) and \(z=v\) gives
\[
  \|u\times v\|^2=\|u\|^2\|v\|^2-(u\cdot v)^2.
\]
This is the area formula for the parallelogram spanned by \(u\) and \(v\).  It also shows that Cauchy--Schwarz is equivalent to the nonnegativity of a squared area.

\section{Why tensors enter naturally}

The cross product is special to three dimensions, but the ideas behind it are not.  The determinant, alternating forms, metrics, and contractions are tensorial constructions.  A tensor is not a multidimensional array by itself; the array is a coordinate representation of a multilinear object.

For example, the metric tensor lets us convert vectors to covectors.  The volume form lets us express oriented volume.  In three-dimensional Euclidean space, the cross product is essentially produced by combining the metric with the volume form.  This explains why the cross product behaves so specially and why it does not generalize naively to every dimension.

\section{Coordinate-free moral}

The note's vector identities are training for coordinate-free thinking.  Coordinates are useful for checking signs, but the real objects are projections, oriented areas, volumes, and multilinear maps.  Once this is understood, tensor notation stops looking like decoration: it is a way to write formulas that remain meaningful after a change of basis.

\section{The hidden structure}

The vector triple product is a projection identity; the scalar triple product is a volume form; the cross product is a Hodge-star translation of an oriented area element.  These three statements give the structural content of the note.

\begin{proposition}[Cross product from wedge product and Hodge star]
In an oriented Euclidean three-dimensional vector space,
\[
  a\times b=\star(a\wedge b).
\]
Consequently,
\[
  c\cdot(a\times b)=\operatorname{vol}(c,a,b),
\]
where \(\operatorname{vol}\) is the oriented volume form.
\end{proposition}

\begin{proof}
The Hodge star is defined so that pairing a vector \(c\) with \(\star(a\wedge b)\) gives the oriented volume of \(c,a,b\).  In the standard oriented orthonormal basis, this is exactly the usual determinant formula for \(c\cdot(a\times b)\).  Therefore \(\star(a\wedge b)\) has the same defining inner products as \(a\times b\).
\end{proof}

This is the right high-level stopping point for the present note.  General basis-change rules and tensor representations belong more naturally to the following linear-map note.

\section{Levi-Civita symbol as a tensor}

The vector identities in this note become systematic with the Levi-Civita symbol.  In an oriented orthonormal basis of $\mathbb R^3$,
\[
  \varepsilon_{ijk}=
  \begin{cases}
  1,&(i,j,k)\text{ even permutation of }(1,2,3),\\
  -1,&(i,j,k)\text{ odd permutation},\\
  0,&\text{otherwise}.
  \end{cases}
\]
Then
\[
  (u\times v)_i=\sum_{j,k}\varepsilon_{ijk}u_jv_k.
\]
The scalar triple product is
\[
  u\cdot(v\times w)=\sum_{i,j,k}\varepsilon_{ijk}u_iv_jw_k.
\]
This notation reveals the tensorial structure behind determinant formulas.

\section{Contraction identities}

Many vector identities reduce to contractions of $\varepsilon_{ijk}$.  The key identity is
\[
  \sum_k \varepsilon_{ijk}\varepsilon_{mnk}
  =\delta_{im}\delta_{jn}-\delta_{in}\delta_{jm}.
\]
Using this, one obtains the vector triple product formula
\[
  u\times(v\times w)=v(u\cdot w)-w(u\cdot v).
\]
What looks like a clever mnemonic is actually an index contraction identity.

\section{Coordinate-free meaning of vector identities}

A tensor is not its array of components.  Components change when the basis changes, while the tensorial object remains.  A bilinear form, for instance, is a map
\[
  B:V\times V\to k.
\]
Its matrix depends on a basis, but the operation $B(u,v)$ does not.

The elementary formulas of dot product, cross product, and triple product are therefore best read as components of coordinate-free objects: metrics, alternating forms, and contractions.

\section{Returning to the original vector identities}

The original note's identities are not separate facts to memorize.  Dot products come from a metric, cross products come from wedge product plus Hodge star, and triple products come from the volume form.  The tensor notation gives a uniform way to prove and remember all of them.

\section{From vector identities to index-free proofs}

The vector triple product
\[
  a\times(b\times c)=b(a\cdot c)-c(a\cdot b)
\]
can be derived by index contraction, but it also has a geometric meaning.  The vector $b\times c$ is perpendicular to the plane spanned by $b,c$, so $a\times(b\times c)$ lies in that plane.  The right-hand side expresses it as the unique linear combination of $b$ and $c$ with coefficients determined by inner products with $a$.

In exterior algebra, this identity is a special case of relationships between wedge product, contraction, and Hodge star.  If $\iota_a$ denotes contraction by $a$, then expressions of the form
\[
  \iota_a(b\wedge c)
\]
encode projections of $a$ against an oriented area element.  The elementary formula is therefore a low-dimensional coordinate shadow of a general calculus of multivectors.

\section{Why tensor notation matters}

Einstein notation is not merely shorter writing.  It separates free indices from contracted indices.  Free indices describe the type of the output; repeated indices describe summation.  For example,
\[
  (a\times b)_i=\varepsilon_{ijk}a_jb_k
\]
has one free index $i$, so the result is a vector.  The repeated indices $j,k$ are summed.

This discipline prevents type errors.  A scalar has no free indices, a vector has one, a matrix has two, and higher tensors have more.  The original vector computations become much easier to audit once every expression has a tensor type.
